\documentclass[11pt,a4paper]{article}
% Estilo comum dos PDFs individuais da Sprint 1.
\usepackage[utf8]{inputenc}
\usepackage[T1]{fontenc}
\usepackage[brazilian]{babel}
\usepackage[a4paper,margin=2.35cm]{geometry}
\usepackage{xcolor}
\usepackage{hyperref}
\usepackage{enumitem}
\usepackage{listings}
\usepackage{parskip}
\usepackage{needspace}
\definecolor{azul}{RGB}{20,63,112}
\hypersetup{colorlinks=true,linkcolor=azul,urlcolor=azul}
\setlist[itemize]{leftmargin=*,itemsep=3pt}
\lstset{basicstyle=\ttfamily\small,frame=single,breaklines=true,columns=fullflexible}
\sloppy
\newcommand{\codigo}[1]{\texttt{#1}}
\newcommand{\titulo}[2]{\begin{center}{\LARGE\bfseries #1\par}\vspace{4pt}{\large #2\par}\vspace{8pt}\rule{0.8\textwidth}{0.4pt}\par\vspace{5pt}{\small Matheus Moreli Durão Silva -- Matrícula 12621CBS217\par Cibersegurança - Bacharelado - Noturno\par}\end{center}}
\newcommand{\cardfonte}[1]{\par\smallskip\noindent\textit{Referência: }\href{#1}{Trello - enunciado e checklist do card}.}

\begin{document}\titulo{S1\_14 - Persistência e JSON}{Guardar vulnerabilidades fora da memória}
\section*{Entregável correto}
O exercício pede um programa de registro de vulnerabilidades em JSON. Persistir significa que os dados continuam disponíveis depois que o programa termina. Sem arquivo, uma lista em Python desaparece quando a execução acaba.
\section*{JSON e serialização}
JSON é um formato textual estruturado, parecido com dicionários e listas: chaves e textos usam aspas duplas. \codigo{json.dump(dados, arquivo)} serializa objetos simples do Python para arquivo; \codigo{json.load(arquivo)} faz o caminho contrário. Deve-se abrir para escrita com \codigo{"w"}, leitura com \codigo{"r"} ou acréscimo de texto com \codigo{"a"}. CSV é útil para tabelas simples; JSON é mais adequado aqui porque uma vulnerabilidade pode ter vários campos.
\begin{lstlisting}
import json
vulnerabilidade = {"descricao": "senha padrao", "severidade": "Alta"}
with open("vulnerabilidades.json", "w", encoding="utf-8") as arquivo:
    json.dump([vulnerabilidade], arquivo, ensure_ascii=False, indent=4)
\end{lstlisting}
\codigo{with} fecha o arquivo automaticamente, mesmo se algo der errado. \codigo{ensure\_ascii=False} preserva acentos e \codigo{indent=4} deixa o arquivo legível.
\section*{Aplicação no projeto}
O \codigo{app.py} usa JSON para o inventário completo. Antes de salvar, os dados precisam ser compostos somente por tipos suportados, como dicionário, lista, texto e número. Por isso, valores de enumeração são convertidos para uma forma simples.
\section*{Checklist demonstrado}
Persistência, leitura/escrita de texto, modos \codigo{r/w/a}, CSV, JSON, serialização/desserialização e erros de arquivo foram estudados. O teste do entregável é executar, abrir o JSON e conferir que a vulnerabilidade está registrada em estrutura válida.
\cardfonte{https://trello.com/c/xdRh1TRi}
\end{document}
